
This paper asked: when RL agents learn to generate code, how much of each generated token actually matters? The answer provided by this mechanism validation study is: \textbf{those that execute.}

By decomposing FGO into three testable components and verifying each with explicit falsification criteria, the following is established:

\begin{enumerate}
\item Trace collection is reliable (100\% capture rate)
\item Gradient exclusion is correct (zero gradient for masked tokens in all checks)
\item The mechanism improves performance (+10\% pass@1 in simulation, 1.78x signal concentration)
\end{enumerate}

Token-level classification achieves 81\% F1, sufficient to demonstrate the mechanism but leaving room for improvement through AST-based span mapping.

The mechanism validation approach demonstrated here---decomposing techniques into independently testable components with falsification criteria---provides a template for rigorous evaluation of credit assignment methods in code RL. This approach enables precise failure localization and avoids conflation of multiple components that has characterized prior work.

Claims regarding convergence speed and content-versus-granularity comparisons remain unverified and require full training runs. The validated mechanism components (trace collector, token classifier, masked PPO loss) are provided for future research.
